\chapter*{Tóm lược nội dung các chương}

\paragraph{Chương 1 — Mở đầu.} Trình bày động lực của việc dùng mô hình ngôn
ngữ nhỏ (SLM) làm web agent, phát biểu bài toán tự động sinh dữ liệu huấn
luyện chất lượng cao từ một website không có nhãn, và nêu mục tiêu, đóng góp,
phạm vi của khóa luận.

\paragraph{Chương 2 — Kiến thức nền tảng và công trình liên quan.} Giới
thiệu bài toán web agent và benchmark WebArena, sau đó khảo sát các hướng:
web agent dựa trên LLM, khám phá skill và workflow, sinh dữ liệu huấn luyện
tự động, và SLM trong vai trò agent; cuối chương định vị khóa luận so với
từng hướng.

\paragraph{Chương 3 — Phương pháp đề xuất.} Trình bày HSGE, một pipeline
khám phá website dựa trên graph: biểu diễn website bằng graph hai tầng, khám
phá có định hướng với các low-level task do LLM đề xuất, phát hiện và kiểm
chứng skill, sinh semantic task bằng cách ghép chuỗi các operator, tinh lọc
dữ liệu, và fine-tune một SLM bằng LoRA.

\paragraph{Chương 4 — Thực nghiệm.} Mô tả môi trường thực nghiệm trên các
website WebArena tự triển khai, báo cáo kết quả thu thập dữ liệu và
fine-tuning, và đánh giá SLM đã fine-tune so với mô hình gốc trên bộ task
chính thức của bốn website: Shopping, Shopping Admin, GitLab và Reddit.

\paragraph{Chương 5 — Thảo luận.} Diễn giải kết quả và đúc kết kinh nghiệm
kỹ thuật thu được khi vận hành LLM agent trên website thật thành các bài học
và best practice có thể tái sử dụng, cùng các nguy cơ đối với tính hợp lệ
(threats to validity).

\paragraph{Chương 6 — Kết luận.} Tổng kết sản phẩm và đóng góp, đối chiếu
kết quả đạt được với mục tiêu ban đầu, và đề xuất hướng phát triển.
